% ---
% id: reductions/rem:intertwining
% kind: rem
% title: "Classical maps before a squasher, and relabelling after one"
% summary: >-
%   A noise or coarse-graining map changes the physical family a squasher is
%   built from, so it acts on the physical POVM before the first link. A
%   stochastic relabelling of both sides by the private setting after a
%   link
%   keeps the channel, the sectors, every pull-back and every certificate.
% prominence: standard
% status: accepted
% uses: [lem:linear, def:joint-family]
% mentions: [def:link, def:chain, ass:tsm, rem:master-scope, lem:master, req:post-processing-order]
% assumes: [A-TSM]
% refs: []
% provenance: newly-derived
% lean:
% ---

Lemma~\ref{lem:linear} post-processes both sides of a relation at once
(part (i)) and folds a private variable into the outcome (part (ii)). This
remark says where in a chain of squashers each of the two may be used.

\begin{remark}\label{rem:intertwining}
\emph{A private variable inside the alphabet.} In
Lemma~\ref{lem:linear}\,(ii) the private variable never enters $\Qsys$: the
target space is unchanged, no block is opened, and no constraint is produced,
the weights $q_t$ being carried by the target effects themselves. A variable
the measuring party draws privately, such as a probe setting or an internal coin of known
bias, is therefore squashed into the outcome alphabet at no cost \emph{by (ii)
alone}. That is the treatment the output-state test-statistics model of
Assumption~\ref{ass:tsm} requires, and the reason
Definition~\ref{def:joint-family} builds the joint alphabet before squashing
rather than after. The no-cost statement concerns (ii) alone and need not
survive a squasher composed with it: once (ii) is composed with a flag-shaped
squasher, every value of the private variable opens flag sectors of its own.
A variable the squasher's channel does not depend on, such as a test-round designation drawn against the setting, is
then carried more cheaply \emph{after} the link, by (i), than inside the
alphabet before it (the relabelling \eqref{eq:relabel}); only a
variable the channel depends on, as a probe setting does through the
conditional families, has to be in the alphabet before the first link. The
branching rule of Lemma~\ref{lem:master}, by contrast, files its classical
variable in $\Qsys$ and does cost something (Remark~\ref{rem:master-scope}).

Part (i) applies a post-processing to \emph{both} sides at once, and two uses
of it differ: one is refused across a link and the other is not.

\emph{A noise or coarse-graining map changes the family a squasher is built
from.} A \tref{dark-count channel} or a coarse-graining applied to the
physical POVM produces a different physical family, and a squasher
constructed for the original one is not a squasher to the same target for it:
$M\dsq_{\Lambda}F$ gives $\Lambda^{\dagger}(F_v)=M_v$, which is not $(PM)_v$
in general. That is why such a map is applied to the physical POVM
\emph{before} the first squashing link (in the \tref{weight route} of the
flag-state submodule, the descended dark-count matrix acts before the joint
family of Definition~\ref{def:joint-family} is formed). Part (i) does move a
stochastic $P$ across a fixed squasher, but only together with the target:
from $M\dsq_{\Lambda}F$ it gives $PM\dsq_{\Lambda}PF$, with the same channel
and the target $PF$ in place of $F$, whose elements, and with them its
\tref[scalar value]{scalar values}, its class sums and the entries built from
them, are not the ones the chain was built and certified for. Where the map
acts on a finer alphabet than the squasher's --- a dark-count channel $P$ on
raw click patterns, applied before the \tref{declared coarse-graining} $K$
that produces the
squasher's family $KM_{\mathrm{raw}}$ from the raw click-pattern POVM
$M_{\mathrm{raw}}$ --- moving it across the squasher asks for a stochastic
$P'$ on the coarse alphabet with $P'KM_{\mathrm{raw}}=KPM_{\mathrm{raw}}$, a
linear feasibility problem in $P'$ with no solution in general. For two
detectors whose single-click patterns $K$ merges into one letter, with
dark-count probabilities $d_1\ne d_2$ acting independently, the double-click
element of $KPM_{\mathrm{raw}}$ carries the element of the pattern in which
only detector~$1$ clicked with the coefficient $d_2$ and that of the pattern
in which only detector~$2$ clicked with the coefficient $d_1$, while every
element of $P'KM_{\mathrm{raw}}$ carries the two with one common
coefficient; so no $P'$ exists when the four elements of $M_{\mathrm{raw}}$
are linearly independent. This contract does neither: the order is fixed
(Requirement~\ref{req:post-processing-order}).

\emph{A stochastic relabelling applied to both sides after a link is part (i)
itself, and it is sound.} Let $M\dsq_{\Lambda}F$ be a link's relation over
the joint alphabet $\Vset$ of Definition~\ref{def:joint-family}, and let
$P(f\,|\,t)$ be a stochastic matrix from the settings $\Tset$ to a finite
alphabet of values $f$, depending on the private component $t=t(v)$ of the
outcome only and never on the device outcome $y$, since otherwise $f$ would
be inferred from an outcome. Then, over the support of $P$,
\begin{equation}\label{eq:relabel}
  \big\{P(f\,|\,t(v))\,M_v\big\}_{(v,f)}\ \dsq_{\Lambda}\
  \big\{P(f\,|\,t(v))\,F_v\big\}_{(v,f)}
\end{equation}
by Lemma~\ref{lem:linear}\,(i) for the stochastic matrix from $\Vset$ to
$\Vset\times\{f\}$ with entries $\delta_{vv'}P(f\,|\,t(v'))$ --- every column
sums to $\sum_fP(f|t(v'))=1$ --- with the \emph{same} channel. Nothing of the
link changes: $\Qsys$, its \tref{sector map} and its dimension are unchanged;
every element $P(f|t(v))F_v$ is a scalar multiple of an existing one, so
block diagonality holds by construction; every carried operator has the
same pull-back; and every entry's certificate is untouched. On a
\tref{scalar sector} the value becomes $P(f|t(v))\phi_s(v)$, and no sector is indexed by
$f$. Such a relabelling is therefore not a link in the sense of
Definition~\ref{def:link}: a link composes a new channel, while a relabelling
changes the alphabet at both ends of an existing relation. This contract
declares no kind, no register and no position in a chain
(Definition~\ref{def:chain}) for it: it is the \emph{consumer's} to apply,
after the chain, by the same lemma, which the module that compiles the
protocol may use directly.
\end{remark}
